\documentclass[11pt,a4paper]{article}
\usepackage[T1]{fontenc}
\usepackage[utf8]{inputenc}
\usepackage{lmodern}
\usepackage[margin=1in]{geometry}
\usepackage{amsmath,amssymb,amsthm}
\usepackage{booktabs}
\usepackage{array}
\usepackage[hidelinks]{hyperref}

\newtheorem{theorem}{Theorem}
\newtheorem{conjecture}{Conjecture}
\theoremstyle{remark}
\newtheorem{remark}{Remark}

\newcommand{\lean}[1]{\texttt{#1}}
\newcommand{\N}{\mathbb{N}}

\title{Machine-checked cases of Je\'smanowicz' conjecture in Lean~4,\\
with formal statements of six further open conjectures}
\author{Dr.\ Moises Romero, Ph.D., SPIE CNG}
\date{25 November 2026}

\sloppy
\begin{document}
\maketitle

\begin{abstract}
We report a Lean~4 / mathlib formalization of Je\'smanowicz' conjecture on
Pythagorean triples together with machine-checked proofs of several of its
classical cases: the triples $(3,4,5)$ and $(5,12,13)$ for every multiplier
$n \ge 1$; the infinite family $(4k^2-1,\,4k,\,4k^2+1)$ and the triple
$(7,24,25)$ for $n=1$; part of the family with $c=b+1$; and a general
restriction on the exponents when $n>1$. All of these are known theorems;
the contribution is their formal verification, by elementary arguments, with no
unproved assumptions. \textbf{The conjecture itself remains open and is not
proved here.} The same repository contains formal statements, with partial
results, of six further open conjectures; none of them is resolved either.
\end{abstract}

\section{Introduction}

A triple $(a,b,c)$ of positive integers is a \emph{primitive Pythagorean
triple} if $a^2+b^2=c^2$ and $\gcd(a,b)=1$.

\begin{conjecture}[Je\'smanowicz, 1955/56]\label{conj}
Let $(a,b,c)$ be a primitive Pythagorean triple and $n\ge 1$ an integer. If
$x,y,z$ are positive integers with
\begin{equation}\label{eq:main}
(an)^x+(bn)^y=(cn)^z ,
\end{equation}
then $x=y=z=2$.
\end{conjecture}

The conjecture originates in a paper of Je\'smanowicz~\cite{Jes} that followed
Sierpi\'nski's treatment of $3^x+4^y=5^z$~\cite{Sier}. Both appeared in a
volume dated 1955/56, and the literature cites either year. We did not consult
the Polish originals, so we cannot say whether the conjecture was first posed
only for $n=1$ or for general $n$; sources differ on this point. For the state
of the art see the surveys \cite{LSS,SDCT}. As far as a brief literature search
in November 2026 could establish, the conjecture is open.

This note describes a formalization in the Lean~4 proof assistant with the
mathlib library. Every theorem stated in Section~\ref{sec:results} has been
checked by the Lean kernel.

\section{Formal statement}

Working over the natural numbers, the formalization defines the following.
\begin{itemize}
\item \lean{IsPrimitivePythagorean}~$a\ b\ c$ means $0<a$, $0<b$, $0<c$,
$\gcd(a,b)=1$ and $a^2+b^2=c^2$;
\item \lean{IsSolution}~$a\ b\ c\ n\ x\ y\ z$ means $0<x$, $0<y$, $0<z$ and
$(an)^x+(bn)^y=(cn)^z$;
\item \lean{JesmanowiczFor}~$a\ b\ c\ n$ means that every $x,y,z$ with
\lean{IsSolution}~$a\ b\ c\ n\ x\ y\ z$ satisfy $x=y=z=2$;
\end{itemize}
It states Conjecture~\ref{conj} as the proposition
\lean{JesmanowiczConjecture}: for all $a,b,c,n$, if
\lean{IsPrimitivePythagorean}~$a\ b\ c$ and $0<n$ then
\lean{JesmanowiczFor}~$a\ b\ c\ n$.
This proposition is only \emph{stated} in the repository. No proof of it is
given or claimed.

\section{Machine-checked results}\label{sec:results}

Throughout, $x,y,z$ are positive integers. The name in typewriter type after
each statement is the Lean theorem.

\begin{theorem}[all multipliers, two triples]\label{thm:alln}
For every $n\ge 1$:
\begin{enumerate}
\item $(3n)^x+(4n)^y=(5n)^z$ implies $x=y=z=2$ \\ \textup{(\lean{jesmanowiczFor\_3\_4\_5})};
\item $(5n)^x+(12n)^y=(13n)^z$ implies $x=y=z=2$ \\ \textup{(\lean{jesmanowiczFor\_5\_12\_13})}.
\end{enumerate}
\end{theorem}

\begin{theorem}[an infinite family, $n=1$]\label{thm:lu}
For every $k\ge 1$,
\[
(4k^2-1)^x+(4k)^y=(4k^2+1)^z \quad\text{implies}\quad x=y=z=2
\]
\textup{(\lean{jesmanowiczFor\_lu})}. Moreover $7^x+24^y=25^z$ implies $x=y=z=2$
\textup{(\lean{jesmanowicz\_7\_24\_25})}.
\end{theorem}

\begin{theorem}[exponent restriction for $n>1$]\label{thm:red}
Let $(a,b,c)$ be a primitive Pythagorean triple and $n>1$. Every solution of
\eqref{eq:main} satisfies $x=y=z=2$, or $x<z<y$, or $y<z<x$
\\ \textup{(\lean{IsSolution.strict\_cases\_of\_one\_lt})}.
\end{theorem}

\begin{theorem}[the family $c=b+1$, $n=1$, partial]\label{thm:dem}
Let $k\ge 1$ and $(a,b,c)=(2k+1,\ 2k(k+1),\ 2k(k+1)+1)$.
\begin{enumerate}
\item Every solution of $a^x+b^y=c^z$ with $y\ge 2$ has $x=y=z=2$
\\ \textup{(\lean{demjanenko\_eq\_two\_of\_two\_le\_y})}.
\item A solution with $y=1$ forces $x\equiv 0 \pmod 4$, $z\equiv 1\pmod 4$, and
every prime factor of $2k+1$ to be $\equiv 1 \pmod 8$ \textup{(\lean{dem\_y\_one})}.
\item Consequently \lean{JesmanowiczFor}~$a\ b\ c\ 1$ holds whenever $4\nmid k$
\\ \textup{(\lean{jesmanowiczFor\_demjanenko\_of\_not\_four\_dvd})}.
\end{enumerate}
\end{theorem}

\begin{remark}
Theorem~\ref{thm:dem} does \emph{not} establish the whole family. The case
$4\mid k$ with all prime factors of $2k+1$ congruent to $1$ modulo $8$ is
reduced in the formalization to an explicitly unproved statement; the smallest
triple not covered is $(17,144,145)$. In the literature the full family is
attributed to Dem'janenko~\cite{Dem}, whose paper we were unable to obtain.
\end{remark}

\subsection*{Proof ideas}

All arguments are elementary. A size comparison shows that a solution with
$x\le z$ and $y\le z$ is $(2,2,2)$. For Theorem~\ref{thm:alln}(1) the remaining
cases are split by the position of $z$. If $z<\min(x,y)$, cancelling $n^z$
shows $n\mid 5^z$ and the case reduces to Sierpi\'nski's theorem. If
$y\le z<x$ or $x\le z<y$, cancelling shows that $n$ is a power of $2$ or of
$3$ and that the cofactor equals $1$, leaving divisibilities of the form
$3^x\mid 5^z-c^z$ or $4^y\mid 5^z-c^z$ with small $c$; these contradict
lifting-the-exponent bounds. For $(5,12,13)$ the multiplier need not be a prime
power, and the proof instead classifies the cofactor $12^y/n^{z-y}$, which is
coprime to $n$, as one of $1,3^y,4^y,12^y$. Theorem~\ref{thm:lu} uses the
congruence $u^2 \mid (1+u)^m-1-mu$ with $u=4k^2$ to force $x$ and $z$ even and
$y\ge 2$, followed by a difference-of-squares factorization in which one
factor must equal~$2$. Theorem~\ref{thm:red} rests on the lemma
$b^{x+1}\nmid c^x-a^x$ for primitive triples.

\subsection*{Attribution}

None of the results above is new. Theorem~\ref{thm:alln}(1) for $n=1$ is due to
Sierpi\'nski~\cite{Sier}; the cases $(5,12,13)$ and $(7,24,25)$ for $n=1$ to
Je\'smanowicz~\cite{Jes}; the family of Theorem~\ref{thm:lu} to Lu~\cite{Lu}.
The all-$n$ statements of Theorem~\ref{thm:alln} are commonly attributed to Deng
and Cohen~\cite{DC}; we did not verify this against the paper, and for
$(5,12,13)$ the attribution should be treated as unconfirmed.
Theorem~\ref{thm:red} is attributed in \cite{note} to Deng~\cite{Deng} and to
Yang and Fu~\cite{YF}. Theorem~\ref{thm:dem}(3) is attributed in \cite{SDCT} to
Podyspanin. Bibliographic data for \cite{Lu,Dem,Deng,YF} were taken from
secondary sources.

The formal proofs were written from scratch and were not transcribed from the
original papers, so they may differ from the published arguments. One exception
should be recorded: the argument for Theorem~\ref{thm:dem}(1) is not independent
of Chapter~1 of Miyazaki's thesis~\cite{Miy}, which was consulted while it was
being completed.

\section{Verification}

The development uses Lean~4 (version 4.36.0) and mathlib. The command
\lean{lake build} succeeds on the whole library. The files contain no
\lean{sorry}, no added axioms and no \lean{native\_decide}. Running
\lean{\#print axioms} on each theorem cited above reports at most the three
standard axioms \lean{propext}, \lean{Classical.choice} and \lean{Quot.sound}.

\section{Other formal statements in the repository}

The repository also contains formal statements of six further open conjectures,
each with machine-checked partial results (Table~\ref{tab:other}). In every case
the conjecture is only stated. The partial results are special cases,
reformulations, or implications between open statements.

\begin{table}[ht]
\centering
\small
\begin{tabular}{@{}>{\raggedright}p{0.30\textwidth}>{\raggedright\arraybackslash}p{0.64\textwidth}@{}}
\toprule
Conjecture & Machine-checked partial results \\
\midrule
Elliott--Halberstam (von Mangoldt form) &
Trivial bounds; monotonicity in the level of distribution; reductions to
natural exponents, large $x$, and levels tending to $1$. \\
Hot spots (Rauch) &
Zero mean and sign change of Neumann eigenfunctions; the interval in full. \\
Goldbach (von Mangoldt-weighted) &
Positivity criteria for the weighted counts; explicit bound on the prime-power
contribution; equivalent forms; verification for even $N\le 10^4$. \\
Legendre (Chebyshev $\theta$, $\psi$) &
Reformulations; Oppermann $\Rightarrow$ Legendre; Andrica $\Rightarrow$
Legendre; verification for $n\le 20000$. \\
P\'olya's eigenvalue conjecture &
The interval in full; lattice-point bound giving the inequality for the
rectangle's explicit model spectrum. \\
Hardy--Littlewood prime $k$-tuples (von Mangoldt form) &
The singular series converges and is positive exactly for admissible tuples;
the quantitative form implies the qualitative form, which implies the twin
prime conjecture. \\
\bottomrule
\end{tabular}
\caption{Further formal statements. None of these conjectures is resolved.}
\label{tab:other}
\end{table}

\section{Limitations}

\begin{itemize}
\item \textbf{Nothing open is solved.} Every conjecture in the repository,
including Conjecture~\ref{conj}, remains unproved.
\item \textbf{Numerical checks are evidence only.} The Goldbach and Legendre
verifications cover ranges far below those in the published literature.
\item \textbf{Statement fidelity.} Mathlib has no Sobolev spaces. The hot spots
and P\'olya statements therefore use weak or classical eigenfunction
definitions whose agreement with the standard ones is argued informally, and
the rectangle case of P\'olya's inequality is proved only for the explicit
model spectrum, without a completeness theorem.
\item \textbf{Literature checks were brief.} Open status and attributions rest
on a small number of searches and largely on secondary sources. In particular,
unrefereed 2026 preprints claim substantial parts of the planar hot spots
conjecture and of P\'olya's conjecture for simply connected planar domains; we
did not verify them.
\end{itemize}

\section*{Availability}

Source code, per-conjecture issue logs and literature notes:
\url{https://github.com/boogiepazzo/romero-conjecture}.

\section*{Acknowledgement}

The Lean proofs and the literature notes were developed with the assistance of
an AI coding agent (Claude Code, Anthropic). All formal statements cited here
are checked by the Lean kernel; the informal attributions and literature
remarks are not machine-checked and carry the uncertainties stated above.

\begin{thebibliography}{99}

\bibitem{Jes} L. Je\'smanowicz, Several remarks on Pythagorean numbers (Polish),
\emph{Wiadom. Mat.} \textbf{1} (1955/56), 196--202.

\bibitem{Sier} W. Sierpi\'nski, On the equation $3^x+4^y=5^z$ (Polish),
\emph{Wiadom. Mat.} \textbf{1} (1955/56), 194--195.

\bibitem{Lu} W. T. Lu, On the Pythagorean numbers $4n^2-1$, $4n$ and $4n^2+1$
(Chinese), \emph{Acta Sci. Natur. Univ. Szechuan} \textbf{2} (1959), 39--42.

\bibitem{Dem} V. A. Dem'janenko, On Je\'smanowicz' problem for Pythagorean
numbers (Russian), \emph{Izv. Vys\v{s}. U\v{c}ebn. Zaved. Mat.} \textbf{48}
(1965), 52--56.

\bibitem{DC} M.-J. Deng and G. L. Cohen, On the conjecture of Je\'smanowicz
concerning Pythagorean triples, \emph{Bull. Austral. Math. Soc.} \textbf{57}
(1998), 515--524.

\bibitem{Deng} M.-J. Deng, A note on the Diophantine equation
$(na)^x+(nb)^y=(nc)^z$, \emph{Bull. Aust. Math. Soc.} \textbf{89} (2014),
316--321.

\bibitem{YF} H. Yang and R.-Q. Fu, A note on Je\'smanowicz' conjecture
concerning primitive Pythagorean triples, \emph{J. Number Theory} \textbf{156}
(2015), 183--194.

\bibitem{Miy} T. Miyazaki, doctoral thesis, Tokyo Metropolitan University,
\url{https://tokyo-metro-u.repo.nii.ac.jp/record/2119/files/10432-001.pdf}.

\bibitem{LSS} M. Le, R. Scott and R. Styer, A survey on the ternary purely
exponential Diophantine equation $a^x+b^y=c^z$, arXiv:1808.06557.

\bibitem{SDCT} G. Soydan, M. Demirci, I. N. Cang\"ul and A. Togb\'e, On the
conjecture of Je\'smanowicz, arXiv:1706.05480.

\bibitem{note} A note on Je\'smanowicz' conjecture for non-primitive
Pythagorean triples, arXiv:2102.10921.

\end{thebibliography}

\end{document}
